\subsection{Universal fiber bundles}

Let G be a topological group, let B and B' be topological spaces, and let \((\mathrm{E}, p)\) and \((\mathrm{E}', p')\) be principal fiber bundles with structure group G and base spaces B and B', respectively.

Let U be an open subset of B, and let \(f'\colon \overline{p}^{-1}(U) \to E'\) be a continuous map that is compatible with the operations of G in \(\overline{p}^{-1}(U)\) and \(E'\), respectively. There then exists a unique continuous map \(f\colon U \to B'\) such that \(f \circ p_{U} = p' \circ f'\) and the commutative square

\[
\begin{array}{c} \mathrm {E_ {U}} \xrightarrow {f ^ {\prime}} \mathrm {E^ {\prime}} \\ \Biggl \downarrow p _ {\mathrm{U}} \\ \mathrm{U} \xrightarrow {f} \mathrm {B^ {\prime}} \end{array}
\]

is then Cartesian (I, p. 94, example (FP)).

For every open subset U of B, denote \(\mathcal{F}(\mathrm{U})\) as the set of continuous maps \(g\colon \mathrm{E}_{\mathrm{U}}\to \mathrm{E}'\) compatible with the operations of G in \(\overline{p}^{-1}(\mathrm{U})\) and \(\mathrm{E}'\), respectively. For every pair (U,V) of open sets of B such that \(\mathrm{U}\subset \mathrm{V}\), let \(r_{\mathrm{UV}}\colon \mathcal{F}(\mathrm{V})\to \mathcal{F}(\mathrm{U})\) be the map defined by \(r_{\mathrm{UV}}(g) = g\mid \mathrm{E}_{\mathrm{U}}\). One checks immediately that one has thus defined a sheaf \(\mathcal{F} = ((\mathcal{F}(\mathrm{U})),(r_{\mathrm{UV}}))\) on B. We shall call this sheaf the sheaf on B of morphisms of principal fiber bundles with group G from E to \(\mathrm{E}'\).

PROPOSITION 4. --- If the space B is paracompact and if the space E' is homotopy equivalent to a point, the sheaf on B of morphisms of principal fiber bundles with group G from E to E' is a soft sheaf.

There exists an open covering \((\mathrm{U}_{j})_{j\in\mathrm{J}}\) of B such that, for every \(j\in J\) , the fiber bundle \(E_{U_{j}}\) is trivializable. Since the space B is paracompact, one may assume the covering \((\mathrm{U}_{j})_{j\in\mathrm{J}}\) locally finite (TG, IX, p. 49) and choose a covering \((\mathrm{A}_{j})_{j\in\mathrm{J}}\) of B where, for every \(j\in J\) the set \(A_{j}\) is closed in B and contained in \(U_{j}\) (TG, IX, p. 49, prop. 4 and p. 48, cor. 1).

By I, p. 65, cor. 2 of prop. 6, it suffices, to prove the proposition, to establish the following assertion: let U be an open subset of B such that the principal fiber bundle \((E_{U}, p_{U})\) is trivializable, let A be a closed subset of B contained in U, let V be an open neighborhood of A contained in U, and let f be an element of \(\mathcal{F}(V)\); then there exists an open neighborhood W of A in V and an element \(f'\) of \(\mathcal{F}(U)\) such that \(r_{\mathrm{WU}}(f') = r_{\mathrm{WV}}(f)\). Let us prove this assertion.

Let W be an open subset of B such that \(A \subset W \subset \overline{W} \subset V\). Let \(s: U \to E_{U}\) be a section of \((\mathrm{E}_{\mathrm{U}}, p_{\mathrm{U}})\). Apply Corollary 3 (IV, p. 438) to the space B, to the closed set \(\overline{W}\), to the neighborhood V of \(\overline{W}\), and to the map \(g = f \circ (s\textbar{} V)\) from \(V\) into \(E'\). There therefore exists a continuous map \(\widetilde{g}: B \to E'\) which coincides with g on \(\overline{W}\). Let \(h = \widetilde{g}\textbar{} U\). We have \(h\textbar{} W = g\textbar{} W = f \circ (s\textbar{} W)\).

The map H: U × G → E' defined by H(x, g) = h(x) · g for (x, g) ∈ U × G is continuous and compatible with the operations of G in U × G and E'. Set \(f' = H \circ s^{-1}\) ; it is an element of \(\mathcal{F}(U)\) . For all \(x \in W\) , the maps f and \(f'\) coincide at the point \(s(x)\) , hence at every point of \(\overline{p}^{-1}(x)\) since these are morphisms of principal fiber bundles. The proposition follows.

THEOREM 1. --- Let G be a topological group, let \(B_{u}\) be a topological space, and let \((\mathrm{E}_{\mathrm{u}}, p_{\mathrm{u}})\) be a principal fiber bundle with group G and base \(B_{u}\). We assume that the space \(E_{u}\) is homotopy equivalent to a point.

Let B be a paracompact topological space.

\begin{enumerate}
\def\labelenumi{\alph{enumi})}
\item
  Every principal fiber bundle with group G and base B is isomorphic to a principal fiber bundle of the form \(f^{*}E_{u}\) , where \(f: B \to B_{u}\) is a continuous map.
\item
  Let \(f_{0}\) and \(f_{1}\) be continuous maps from B into \(B_{u}\). For \(f_{0}^{*}E_{u}\) and \(f_{1}^{*}E_{u}\) to be isomorphic principal fiber bundles with group G and base B, it is necessary and sufficient that the maps \(f_{0}\) and \(f_{1}\) be homotopic.
\end{enumerate}

In other words, there exists a map from \([B, B_{u}]\) to \(P(B; G)\) which associates with the homotopy class of a continuous map f from B into \(B_{u}\) the isomorphism class of the principal fiber bundle \(f^{*}E_{u}\). This map is bijective.

Let E be a principal fiber bundle with group G and base B. By Prop. 4, the sheaf \(\mathcal{F}\) on B of morphisms of principal fiber bundles from E into \(\mathrm{E_u}\) is soft. Consequently, \(\mathcal{F}(\mathrm{B})\) is nonempty, whence \(a\).

Let \(f_{0}\) and \(f_{1}\) be continuous maps from B into \(B_{u}\). If the maps \(f_{0}\) and \(f_{1}\) are homotopic, the principal fiber bundles \(f_{0}^{*}E_{u}\) and \(f_{1}^{*}E_{u}\) are isomorphic (IV, p. 445, cor. 2). Let us prove the converse. For \(\alpha\in\{0,1\}\), denote by \(E_{\alpha}\) the \(B_{u}\)-principal fiber bundle \(f_{\alpha}^{*}E_{u}\), \(g_{\alpha} \colon E_{\alpha} \to E_{u}\) the first projection and \(p_{\alpha} \colon E_{\alpha} \to B\) the second projection. Let \(i \colon E_{0} \to E_{1}\) be an isomorphism of principal fiber bundles. Let \(p\) be the map \(p_{0} \times Id_{I} \colon E_{0} \times I \to B \times I\).

Since the space B × I is paracompact (TG, IX, p. 70, prop. 17), the sheaf G on B × I of morphisms of principal fiber bundles from E₀ × I into Eᵤ is soft (IV, p. 446, prop. 4). Set A = B × \{0,1\}, U = B × ([0, ½[∪]½, 1]), and define an element g of G(U) by setting

\[
g (x, t) = \left\{ \begin{array}{l l} g _ {0} (x) & \text { for } (x, t) \in \mathrm{E} _ {0} \times [ 0, \frac {1}{2} [, \\ g _ {1} \circ i (x) & \text { for } (x, t) \in \mathrm{E} _ {0} \times ] \frac {1}{2}, 1 ]. \end{array} \right.
\]

Since the sheaf \(\mathcal{G}\) is soft, there exists an element \(h \in \mathcal{G}(\mathrm{B} \times \mathbf{I})\) and an open neighborhood V of A in U such that \(h \mid \mathrm{V} = g \mid \mathrm{V}\); such an element is a continuous map \(\mathrm{H} \colon \mathrm{E}_0 \times \mathbf{I} \to \mathrm{E}_{u}\) compatible with the operations of G and such that \(\mathrm{H}(x,0) = g_0(x)\), \(\mathrm{H}(x,1) = g_1(i(x))\) for all \(x \in \mathrm{E}_0\). There then exists a map \(h' \colon \mathrm{B} \times \mathbf{I} \to \mathrm{B}_{u}\) such that \(h'(p_0(x),t) = p_{\mathrm{u}}(\mathrm{H}(x,t))\) for \(x \in \mathrm{E}_0\) and \(t \in \mathbf{I}\); this map is continuous and is a homotopy joining \(f_0\) to \(f_1\).

COROLLARY. --- Let G be a topological group, and let B and B' be paracompact topological spaces. Let (E,p) and (E',p') be principal fiber bundles with group G and bases B and B', respectively. Suppose that the spaces E and E' are both homotopy equivalent to a point. The spaces B and B' are homotopy equivalent.

Indeed, there exists a continuous map \(f \colon \mathrm{B} \to \mathrm{B}'\) such that the principal fiber bundle E is isomorphic to the principal fiber bundle \(f^* \mathrm{E}'\) and a continuous map \(g \colon \mathrm{B}' \to \mathrm{B}\) such that the principal fiber bundle \(\mathrm{E}'\) is isomorphic to the principal fiber bundle \(g^* \mathrm{E}\) (Theorem 1, (a)). The principal fiber bundles \((g \circ f)^* \mathrm{E}\) and E are then isomorphic, so the map \(g \circ f\) is homotopic to the map \(\mathrm{Id}_{\mathrm{B}}\) (Theorem 1, (b)). Likewise, the map \(f \circ g\) is homotopic to the map \(\mathrm{Id}_{\mathrm{B}'}\).

Let G be a topological group, let B be a topological space, and let E be a principal G-bundle with base B. Suppose that the space E is homotopy equivalent to a point. One says that the principal bundle E is universal for principal G-bundles with paracompact base, and one says that the space B is a classifying space for G. If two classifying spaces for G are paracompact, they are homotopy equivalent. When a classifying space exists, the study of isomorphism classes of principal G-bundles with paracompact base can be regarded as a problem in homotopy theory.

Example. --- Equipped with the map \(p: R \to S^{1}\), \(t \mapsto e^{2\pi it}\), and with the action of Z by translation, the space R is a principal covering with group Z. The space \(S^{1}\) is thus a classifying space for the group Z.

For every discrete group G, we will construct in the next section a metrizable space that is a classifying space for G.

